\scanpage{33}
\uline{Formulaire de} $L,\ \handwritten{\Lambda}$.

$X$ variété projective lisse connexe polarisée par $\handwritten{\unsure{\xi|\zeta}}_X\handwritten{\in} H^{\handwritten{2}}(X)$, $L_X$ opération de cup-produit associé\handwritten{e,} $\handwritten{\Lambda}_X$ opération de degré $-2$ correspondante, $p_i^X$ les projections de $H^{\handwritten{*}}(X)$ sur les morceaux primitifs $P^i(X)$ pour $0\handwritten{\leq}i\handwritten{\leq}2n$, $Y$ une section hyperplane \deleted{\unsure{d}}lisse de $X$, $\handwritten{\xi}_Y\handwritten{\in}H^{\handwritten{2}}(Y)$ la self-intersection (induite par $\handwritten{\xi}_X$), d'où $L_Y,\handwritten{\Lambda}_Y,p_Y^i$ ; $\phi:Y\handwritten{\to}X$ l'immersion, d'où $\phi^{\handwritten{*}}$ et $\phi_{\handwritten{*}}$.

A) \qquad On a
\[
\begin{aligned}
1)\quad &\operatorname{Im}L_X=\mathop{\handwritten{\bigcap}}_{0\handwritten{\leq}i\handwritten{\leq}n}\operatorname{Ker}p_{\deleted{\unsure{X}}i}^{\deleted{\unsure{i}}X},&
 \operatorname{Ker}L_X&=\mathop{\handwritten{\sum}}_{n\handwritten{\leq}i\handwritten{\leq}2n}\operatorname{Im}p_i^X,\\
2)\quad &\operatorname{Im}\handwritten{\Lambda}_X=\mathop{\handwritten{\bigcap}}_{n\handwritten{\leq}i\handwritten{\leq}2n}\operatorname{Ker}p_i^X&
 \operatorname{Ker}\handwritten{\Lambda}_X&=\mathop{\handwritten{\sum}}_{0\handwritten{\leq}i\handwritten{\leq}n}\operatorname{Im}p_i^X,
\end{aligned}
\]
d'où en particulier les formules
\[
\begin{aligned}
1')\quad &p_i^X L_X=0\quad\text{si}\quad 0\handwritten{\leq}i\handwritten{\leq}n,\quad&
 \deleted{\unsure{p}}L_Xp_i^X&=0\quad\text{si}\quad n\handwritten{\leq}i\handwritten{\leq}2n,\\
2')\quad &\deleted{\unsure{\phi}}p_{\deleted{\unsure{X}}i}^X\handwritten{\Lambda}_X=0\quad\text{si}\quad n\handwritten{\leq}i\handwritten{\leq}2n,\quad&
 \handwritten{\Lambda_X}p_i^X&=0\quad\text{si}\quad 0\handwritten{\leq}i\handwritten{\leq}n,
\end{aligned}
\]
et en particulier
\[
p_n^X L_X=L_Xp_n^X=p_n^X\handwritten{\Lambda}_X=\handwritten{\Lambda}_X p_n^X=0\quad. \tag*{3)}
\]

B) \qquad On a
\[
\handwritten{\Lambda}_X L_X=\mathrm{id}_X-\mathop{\handwritten{\sum}}_{\deleted{0}\handwritten{n}\handwritten{\leq}i\handwritten{\leq}\deleted{n}\handwritten{2n}}p_i^X,\qquad
 L_X\handwritten{\Lambda}_X=\mathrm{id}_X-\mathop{\handwritten{\sum}}_{\deleted{\unsure{n}}\handwritten{0}\handwritten{\leq}i\handwritten{\leq}\deleted{\unsure{2n}}\handwritten{n}}p_i^X, \tag*{4)}
\]
\deleted{d'où compte tenu de $1')$ et $2')$ les formules
\[
\begin{aligned}
4')\quad &\handwritten{\Lambda}_X L_X^2\handwritten{=}
 \deleted{\unsure{L_X\Lambda_X L_X}}\handwritten{\inserted{\left(\mathrm{id}_X-\sum_{0\leq i\leq 2n}p_i^X\right)L_X}},\quad
 \deleted{\text{et }\unsure{L_X\Lambda_X^2=\Lambda_XL_X\Lambda_X}}\handwritten{\inserted{L_X^2\Lambda_X=L_X\left(\mathrm{id}_X-\sum_{0\leq i\leq 2n}p_i^X\right)}},\\
4'')\quad & L_X\handwritten{\Lambda}_X L_X=L_X\handwritten{\Lambda}_X L_X\deleted{\unsure{\sum_{0\leq i\leq n}p_i^X=L_X\sum_{n\leq i\leq 2n}p_i^X}}\quad\handwritten{\text{et formules ``transposées''.}}
\end{aligned}
\]}

C) \qquad On a
\[
\phi_{\handwritten{*}}\handwritten{\Lambda}_Y\phi^{\handwritten{*}}=\mathrm{id}_X-\mathop{\handwritten{\sum}}_{0\handwritten{\leq}i\handwritten{\leq}2n}p_i^X
 \handwritten{\inserted{\overset{\text{par }4)}{=}(\Lambda_X L_X+L_X\Lambda_X)-\mathrm{id}_X+p_n^X}} \tag*{5)}
\]
d'où\handwritten{,} \deleted{compte tenu de $4'')$, $1')$ et $2')$}
\handwritten{\inserted{posant $(6)\ T_X=\mathrm{id}_X+\varphi_*\Lambda_Y\varphi^*$,}}
\deleted{\[
 5')\quad \handwritten{\Lambda}_X L_X^2=\deleted{\unsure{L_X\phi_*\Lambda_Y\phi^*}}\handwritten{\inserted{\varphi_*\Lambda_Y\varphi^* L_X}},\quad
 \handwritten{L_X^2\Lambda_X=L_X\varphi_*\Lambda_Y\varphi^*,}
\]}
\handwritten{\inserted{(donc $\Lambda_X L_X+L_X\Lambda_X=T_X-p_n^X$),}}

d'où compte tenu de $\deleted{\unsure{5'|3}}\handwritten{3\unsure{*|'}})$ la formule
\[
\begin{aligned}
\deleted{\unsure{6}}\handwritten{7})\quad
 L_X^2\handwritten{\Lambda}_X+2L_X\handwritten{\Lambda}_{\deleted{\unsure{Y|X}}X}L_X+
 \handwritten{\Lambda}_{\deleted{\unsure{Y|X}}X}L_X^2
 &=\deleted{L_X(\mathrm{id}_X\unsure{+|-}\phi_*\handwritten{\Lambda}_Y\phi^*)\unsure{+2|+}( \mathrm{id}_X\unsure{+|-}\phi_*\handwritten{\Lambda}_Y\phi^*)L_X}\\
 &\handwritten{\deleted{\unsure{2L_X+L_X}}\,\inserted{L_XT_X+T_XL_X,\quad\text{avec } T_X}}\\
 &\handwritten{\deleted{T_X=\mathrm{id}_X+\varphi_*\Lambda_Y\varphi^*}\quad\inserted{\text{donné par }(6)}}
\end{aligned}
\]

\uline{Corollaire} \quad $(A(X\mathbin{\mathrm{x}}X)$ et $C(Y))$ implique $C(X)$.

\handwritten{On y peut remplacer $A(X\times X)$ par $A^{\unsure{\circ|0}}(X\times X)$ (condition ``des \unsure{dim|dims|dim.} critiques'' $H^{\unsure{(2n)|2n|(2n+1)}-2}\to H^{(2n)+2}$}
% archive sheet number: 32

